\chapter{期末综合}

\Needspace{12\baselineskip}

\section{期末综合素养评价（一）}

\Needspace{8\baselineskip}

\pdfbookmark[2]{原书第93页}{orig-93}

\subsection*{原书第93页}

% p093_s01
\BookQuestion{一、1}
期末综合素养评价（一）。时间：90分钟；满分：100分（知识：98分，卷面：2分）。\par 一、填空。（每空1分，共15分）\par 1. \(7\)千米\(60\)米＝（\hspace{1.8em}）千米；\allowbreak{}\(800\)平方厘米＝（\hspace{1.8em}）平方米。\par
\begin{BookAnswer}
按进率先换算：\[60\text{米}=0.06\text{千米}\]\[7\text{千米}60\text{米}=(7+0.06)\text{千米}=7.06\text{千米}\]\[1\text{平方米}=10000\text{平方厘米}\]\[800\text{平方厘米}=(800\div10000)\text{平方米}=0.08\text{平方米}\]依次填\(7.06,0.08\)，原空后的单位不重复填写。
\end{BookAnswer}

% p093_s02
\BookQuestion{一、2}
2. 据统计，我国森林覆盖面积达到\(\underline{240960000}\)公顷，横线上的数改写成用“亿”作单位的数是（\hspace{1.8em}）公顷，保留两位小数约是（\hspace{1.8em}）公顷。\par
\begin{BookAnswer}
\[240960000\text{公顷}=2.4096\text{亿公顷}\]保留两位小数看千分位9，进1得\(2.41\text{亿公顷}\)。原题两个空后都只印了“公顷”，因此应在空内分别填\(2.4096\text{亿}\)、\(2.41\text{亿}\)，不能只填2.4096和2.41而丢掉“亿”。这里“亿”是数的数量级标记，补入后与空后的“公顷”合成正确单位。
\end{BookAnswer}

% p093_s03
\BookQuestion{一、3}
3. 妈妈收到一个快递包裹，质量是一个三位小数，它的个位上是\(2\)个一，十分位和千分位上都是\(7\)，百分位上是\(0\)，这个包裹的质量是（\hspace{1.8em}）千克，这个小数的计数单位是（\hspace{1.8em}）。\par
\begin{BookAnswer}
个位是\(2\)，十分位是\(7\)，百分位是\(0\)，千分位是\(7\)。包裹重\(2.707\)千克。最低位是千分位，计数单位可以填“千分之一”或\(0.001\)。
\end{BookAnswer}

% p093_s04
\BookQuestion{一、4}
4. 把一张长方形纸对折两次，打开后两条折痕的位置关系可能是（\hspace{1.8em}）。\par
\begin{BookAnswer}
填“互相垂直或互相平行”。横、竖方向各对折一次，两条折痕可以互相垂直。沿同一方向连续对折，两条折痕可以互相平行。
\end{BookAnswer}

% p093_s05
\BookQuestion{一、5}
5. 小明打算做一个等腰三角形的风筝框架，竹条有\(48\,\mathrm{cm}\)和\(24\,\mathrm{cm}\)两种长度，他围成框架的周长是（\hspace{1.8em}）\(\mathrm{cm}\)。\par
\begin{BookAnswer}
原题只给出竹条有48厘米、24厘米两种长度，没有说明每种各有几根，也没有要求两种长度必须都用。等边三角形也属于等腰三角形，不能不加说明地只填120。\par 若选用三根竹条且相应根数足够，按48厘米竹条的根数分类，只有下面四种无序边长组合：\par \textbf{（1）三根都是24厘米。}组成等边三角形，也就是特殊的等腰三角形。\[\begin{aligned}&24\times3\\={}&72\text{(厘米)}\end{aligned}\]这种情况下，周长为72厘米。\par \textbf{（2）两根24厘米、一根48厘米。}因为\(24+24=48\)，两条短边之和等于最长边，不能围成三角形，排除。\par \textbf{（3）一根24厘米、两根48厘米。}两腰取48厘米、底边取24厘米，满足三角形三边关系。\[\begin{aligned}&48\times2+24\\={}&96+24\\={}&120\text{(厘米)}\end{aligned}\]这种情况下，周长为120厘米。\par \textbf{（4）三根都是48厘米。}组成等边三角形。\[\begin{aligned}&48\times3\\={}&144\text{(厘米)}\end{aligned}\]这种情况下，周长为144厘米。\par 因而，若同长竹条的所需根数足够，周长可能为72、120或144厘米；这三种可能没有把同一组边长的不同摆放重复计数。若额外要求“两种长度都要用”，则只有第（3）种，周长才唯一是120厘米。由于原题未给库存和这一要求，原空不宜作为无条件的唯一答案。
\end{BookAnswer}

% p093_s06
\BookQuestion{一、6}
6. 妈妈买了两只鸭，其中黑鸭重\(1.8\)千克，白鸭比黑鸭重\(700\)克，两只鸭共重（\hspace{1.8em}）千克。\par
\begin{BookAnswer}
白鸭比黑鸭重，先统一单位：\[700\text{克}=0.7\text{千克}\]\[\begin{aligned} &1.8+(1.8+0.7)\\ ={}&1.8+2.5\\ ={}&4.3\text{(千克)}\end{aligned}\]填4.3。
\end{BookAnswer}

% p093_s07
\BookQuestion{一、7}
7. 在○里填上“＞”“＜”或“＝”。\par \(125\times8\times4\) ○ \(125\times24\)\par \(480\div3\div2\) ○ \(480\div5\)\par \(625-125-75\) ○ \(625-(125+75)\)\par \(327-98\) ○ \(327-100-2\)\par
\begin{BookAnswer}
先各算两边，再比较：\par （1）\(125\times8\times4=4000\)，\allowbreak{}\(125\times24=3000\)，填\(>\)。\par （2）\(480\div3\div2=80\)，\allowbreak{}\(480\div5=96\)，填\(<\)。\par （3）两边都是\(625-200=425\)，填\(=\)。\par （4）\(327-98=229\)，\allowbreak{}\(327-100-2=225\)，填\(>\)。
\end{BookAnswer}

% p093_s08
\BookQuestion{一、8}
8. 元旦文艺汇演中，跳舞的有\(36\)人，合唱的有\(42\)人，参加这两项演出的总人数最多可能（\hspace{1.8em}）人，最少（\hspace{1.8em}）人。\par
\begin{BookAnswer}
人数最多时，两项演出没有重复人员，\allowbreak{}\(36+42=78\)人；最少时，跳舞的36人都在合唱的42人中，共42人。依次填\(78,42\)。
\end{BookAnswer}

% p093_s09
\BookQuestion{二、1}
二、选择（将正确答案的序号填在括号里）。（每题2分，共16分）\par 1. 一个三角形中最大的角是\(150^\circ\)，那么这个三角形是（\hspace{1.8em}）。\par
\noindent A. 锐角三角形\par
\noindent B. 直角三角形\par
\noindent C. 钝角三角形\par
\begin{BookAnswer}
选C。最大角\(150^\circ>90^\circ\)，有一个钝角，因此是钝角三角形。
\end{BookAnswer}

\Needspace{8\baselineskip}

\pdfbookmark[2]{原书第94页}{orig-94}

\subsection*{原书第94页}

% p094_s01
\BookQuestion{二、2}
2. 下面各数中，一个“0”也不读出的是（\hspace{1.8em}）。\par
\noindent A. 906.00\par
\noindent B. 900.60\par
\noindent C. 900.06\par
\noindent D. 90.080\par
\begin{BookAnswer}
按原数的写法读，小数部分须逐位读：\par A：\(906.00\)，九百零六点零零；\par B：\(900.60\)，九百点六零；\par C：\(900.06\)，九百点零六；\par D：\(90.080\)，九十点零八零。\par 四项都要读出“零”，所以本题按题面无正确选项。若先把\(900.60\)化简为\(900.6\)，再读为“九百点六”，B才符合，但题面没有给出“先化简”的条件。
\end{BookAnswer}

% p094_s02
\BookQuestion{二、3}
3. 关于\(32\times25\times125\)的计算中，下面计算正确的是（\hspace{1.8em}）。\par
\noindent A. \(32\times25\times125=(125\times32)\times(25\times32)\)\par
\noindent B. \(32\times25\times125=125\times8+25\times4\)\par
\noindent C. \(32\times25\times125=(125\times8)+(25\times4)\)\par
\noindent D. \(32\times25\times125=(125\times8)\times(25\times4)\)\par
\begin{BookAnswer}
选D。把\(32\)拆成\(8\times4\)，利用交换律、结合律：\[\begin{aligned} &32\times25\times125\\ ={}&(125\times8)\times(25\times4)\\ ={}&1000\times100\\ ={}&100000\end{aligned}\]A重复乘了32；B、C把应乘的两组数误写成相加。
\end{BookAnswer}

% p094_s03
\BookQuestion{二、4}
4. 两根木棍连在一起，从头到尾一共长\(70\,\mathrm{cm}\)，其中一根长\(50\,\mathrm{cm}\)，中间重合部分长\(5\,\mathrm{cm}\)，另一根木棍长（\hspace{1.8em}）\(\mathrm{cm}\)。\par
\noindent A. 30\par
\noindent B. 45\par
\noindent C. 25\par
\noindent D. 35\par
\begin{BookAnswer}
选C。两根木棍长度相加，比连接后的总长多算了5厘米重合段，所以另一根应长\[\begin{aligned} &70-50+5\\ ={}&20+5\\ ={}&25\text{(厘米)}\end{aligned}\]
\end{BookAnswer}

% p094_s04
\BookQuestion{二、5}
5. 两个完全一样的梯形不可能拼成（\hspace{1.8em}）。\par
\noindent A. 平行四边形\par
\noindent B. 长方形\par
\noindent C. 六边形\par
\noindent D. 三角形\par
\begin{BookAnswer}
选D。两个完全一样的梯形可沿腰或底拼成平行四边形、特殊情况下的长方形，或六边形；不能拼成三角形。梯形的两底长度不相等，不能在拼合时把其余边全部并成三条直边。
\end{BookAnswer}

% p094_s05
\BookQuestion{二、6}
6. 小丽组\(6\)个人的平均身高是\(140\)厘米，小刚组\(4\)个人的平均身高是\(150\)厘米。这两个组的平均身高是（\hspace{1.8em}）。\par
\noindent A. 145厘米\par
\noindent B. 144厘米\par
\noindent C. 148厘米\par
\noindent D. 146厘米\par
\begin{BookAnswer}
选B。\textbf{方法一：移多补少，简便计算}\par 以140厘米为基准，小刚组4人的总身高比4份140厘米多\((150-140)\times4\)厘米，把这部分平均分给全部10人。\[\begin{aligned}&140+(150-140)\times4\div(6+4)\\={}&140+10\times4\div10\\={}&140+40\div10\\={}&140+4\\={}&144\text{(厘米)}\end{aligned}\]\textbf{方法二：先求两组总身高再平均}\par 两组人数不同，应分别用平均身高乘人数，再用身高总和除以总人数。\[\begin{aligned}&(140\times6+150\times4)\div(6+4)\\={}&(840+600)\div10\\={}&1440\div10\\={}&144\text{(厘米)}\end{aligned}\]两个组的平均身高是144厘米。
\end{BookAnswer}

% p094_s06
\BookQuestion{二、7}
7. 右图这个图形是通过（\hspace{1.8em}）运动方式得到的。\par
\noindent A. 平移\par
\noindent B. 旋转\par
\noindent C. 平移或旋转\par
\BookFigure{figures/p094_s06-01.tex}
\begin{BookAnswer}
选B。图中相同的基本图形围着共同的中心改变方向，形成旋转排列；单纯平移不会改变图形朝向。
\end{BookAnswer}

% p094_s07
\BookQuestion{二、8}
8. 超市新春大促销，某品牌酸奶每盒\(4\)元，买\(4\)赠\(1\)，妈妈带了\(36\)元钱，最多能买（\hspace{1.8em}）盒酸奶。\par
\noindent A. 9\par
\noindent B. 10\par
\noindent C. 11\par
\noindent D. 8\par
\begin{BookAnswer}
选C。可付费买\(36\div4=9\)盒，每买4盒赠1盒，9盒中有2组完整的4盒，可获赠2盒；总计\(9+2=11\)盒。余下付费1盒不能再得到赠品。
\end{BookAnswer}

% p094_s08
\BookQuestion{三、1}
三、计算。（共30分）\par 1. 直接写得数。（6分）\par \(14\times60=\quad6.5+7.09=\quad25\times4\div25\times4=\)\par \(7-1.8=\quad70\div1000=\quad10-0.1=\)\par
\begin{BookAnswer}
按原题行列：\par 第一行：\(840,13.59,16\)；\par 第二行：\(5.2,0.07,9.9\)。\par 要点：\(25\times4\div25\times4\)同级运算从左到右，\allowbreak{}\(100\div25\times4=16\)；\allowbreak{}\(70\div1000=0.07\)。小数加减时对齐小数点。
\end{BookAnswer}

% p094_s09
\BookQuestion{三、2}
2. 列竖式计算。（12分）\par \(12.4+2.76=\textcolor{answerred}{15.16}\)\qquad \(8.05-4.357=\textcolor{answerred}{3.693}\)\par \(28\times304=\textcolor{answerred}{8512}\)\qquad \(6270\div30=\textcolor{answerred}{209}\)\par
\begin{BookAnswer}
小数竖式对齐小数点，才能让相同计数单位对齐；末尾补0不改变大小。\(8.050-4.357\)的千分位不够减，先从百分位退1，换成10个千分之一；百分位余下4，不够减5，十分位又是0，须从个位退1，换成10个十分之一，再取其中1个十分之一换成10个百分之一，按位相减：\par \[\begin{array}{r}12.40\\+2.76\\\hline15.16\end{array}\qquad\begin{array}{r}8.050\\-4.357\\\hline3.693\end{array}\]\par 整数乘法按位相乘：\[\begin{array}{r}28\\\times304\\\hline112\\0\phantom0\\84\phantom{00}\\\hline8512\end{array}\]\par 除法竖式：\[\begin{array}{rr}&209\\\cline{2-2}30\big)&6270\\&60\phantom{00}\\\cline{2-2}&270\\&270\\\cline{2-2}&0\end{array}\]\par 62里面最多有2个30，百位商2，余2；落下7得27，十位不够商1，必须商0；再落下0，\allowbreak{}\(270\div30=9\)，个位商9。检验\(30\times209=6270\)。
\end{BookAnswer}

\Needspace{8\baselineskip}

\pdfbookmark[2]{原书第95页}{orig-95}

\subsection*{原书第95页}

% p095_s01
\BookQuestion{三、3（1）}
3. 计算下列各题，能简便计算的要简便计算。（12分）\par \(40\times[175-(75+67)]\)\par
\begin{BookAnswer}
先用减法性质，将\(175\)与\(75\)相减凑整。\[\begin{aligned}&40\times[175-(75+67)]\\={}&40\times[175-75-67]\\={}&40\times[100-67]\\={}&40\times33\\={}&1320\end{aligned}\]
\end{BookAnswer}

% p095_s02
\BookQuestion{三、3（2）}
\(39\times99+99\)\par
\begin{BookAnswer}
方法一：提取共同的\(99\)，再把\(99\)看作\(100-1\)。\[\begin{aligned}&39\times99+99\\={}&(39+1)\times99\\={}&40\times99\\={}&40\times(100-1)\\={}&4000-40\\={}&3960\end{aligned}\]方法二：将\(99\)看成\(100\)少\(1\)，再把\(-39+99\)合并。\[\begin{aligned}&39\times99+99\\={}&39\times(100-1)+99\\={}&3900-39+99\\={}&3900+(99-39)\\={}&3900+60\\={}&3960\end{aligned}\]
\end{BookAnswer}

% p095_s03
\BookQuestion{三、3（3）}
\(245\times36-145\times36\)\par
\begin{BookAnswer}
用乘法分配律合并。\[\begin{aligned} &245\times36-145\times36\\ ={}&(245-145)\times36\\ ={}&100\times36\\ ={}&3600\end{aligned}\]
\end{BookAnswer}

% p095_s04
\BookQuestion{四、1}
四、实践操作。（共8分）\par 1. 在方格图上画一个钝角三角形、一个等腰直角三角形，并分别画出它们一条边上的高。（4分）\par
\BookFigure{figures/p095_s04-01.tex}
\begin{BookAnswer}
钝角三角形可取三个格点，使一个角大于\(90^\circ\)；所画红色钝角三角形从下方顶点向上边作垂线段，垂足与上边重合并标直角。等腰直角三角形取两条等长且互相垂直的格线作直角边；从直角顶点向斜边所作垂线段是对应高。红色“高”的两端均应落在顶点与垂足上。
\BookFigure{figures/p095_s04-02.tex}
\end{BookAnswer}

% p095_s05
\BookQuestion{四、2}
2. 在方格图上画一个底是\(4\,\mathrm{cm}\)、高是\(2\,\mathrm{cm}\)的平行四边形，并把它沿左下角的点逆时针旋转\(90^\circ\)。（4分）\par
\BookFigure{figures/p095_s05-01.tex}
\begin{BookAnswer}
原平行四边形底占4个单位格，高占2个单位格；以标出的左下角点为旋转中心。逆时针旋转\(90^\circ\)时，中心不动，其余顶点到中心的距离保持不变，方向逆时针转直角，再依次连接对应点。红色原图与旋转图已保留，底、高及直角标记与格点对齐。
\BookFigure{figures/p095_s05-02.tex}
\end{BookAnswer}

% p095_s06
\BookQuestion{五、1}
五、解决问题。（共29分）\par 1. 有一批蔬菜要运出基地，已经运走了\(82.6\)吨，运走的比剩下的多\(13.8\)吨，这批蔬菜共几吨？（5分）\par
\begin{BookAnswer}
\textbf{方法一：先求剩下的质量}\par 运走82.6吨，比剩下的多13.8吨，所以剩下82.6减13.8吨。\[\begin{aligned} &82.6-13.8+82.6\\ ={}&68.8+82.6\\ ={}&151.4\text{(吨)}\end{aligned}\]答：这批蔬菜共151.4吨。\par \textbf{方法二：把剩下部分补到82.6吨}\par 两份82.6吨比原总量多13.8吨，最后减去多补的量。\[\begin{aligned} &82.6+82.6-13.8\\ ={}&165.2-13.8\\ ={}&151.4\text{(吨)}\end{aligned}\]答：这批蔬菜共151.4吨。\par
\end{BookAnswer}

% p095_s07
\BookQuestion{五、2}
2. 小林用\(5\)天时间看完一本书，她前\(3\)天平均每天读\(21\)页，后\(2\)天共读\(27\)页，小林平均每天读多少页？（5分）\par
\begin{BookAnswer}
\textbf{方法一：总页数除以总天数}\par 前3天每天21页，后2天共27页，先求全书页数。\[\begin{aligned} &(21\times3+27)\div5\\ ={}&(63+27)\div5\\ ={}&90\div5\\ ={}&18\text{(页)}\end{aligned}\]答：小林平均每天读18页。\par \textbf{方法二：与每天21页比较}\par 后2天如果也每天21页，应读42页，实际少读15页，将15页平均扣给5天。\[\begin{aligned} &21-(21\times2-27)\div5\\ ={}&21-(42-27)\div5\\ ={}&21-15\div5\\ ={}&21-3\\ ={}&18\text{(页)}\end{aligned}\]答：小林平均每天读18页。\par
\end{BookAnswer}

% p095_s08
\BookQuestion{五、3}
3. 爷爷去市场买了一些鸡蛋和一箱纯牛奶，到家后，奶奶问：“你买了多少千克鸡蛋？”爷爷笑着摇头，说记不清了，请你根据下面的信息，帮他算一算这次共买了多少千克鸡蛋。（5分）\par ●信息1：鸡蛋每千克12元；\par ●信息2：一箱纯牛奶36元；\par 另两条信息见下一页。\par
\begin{BookAnswer}
本题信息跨页；100元付款、找回16元等条件在下一页。完整逐行脱式及多解列在本书原书第96页对应的“五、3（续）”。
\end{BookAnswer}

\Needspace{8\baselineskip}

\pdfbookmark[2]{原书第96页}{orig-96}

\subsection*{原书第96页}

% p096_s01
\BookQuestion{五、3（续）}
五、3（续）。\par ●信息3：爷爷付给收银员100元；\par ●信息4：找回16元。\par 结合上一页信息1、2，算这次共买了多少千克鸡蛋。\par
\begin{BookAnswer}
\textbf{方法一：先求鸡蛋总价}\par 实付总价100减16元；再扣掉一箱牛奶36元，最后除以每千克鸡蛋12元。\[\begin{aligned} &(100-16-36)\div12\\ ={}&(84-36)\div12\\ ={}&48\div12\\ ={}&4\text{(千克)}\end{aligned}\]答：这次买了4千克鸡蛋。\par \textbf{方法二：把牛奶价钱换算成鸡蛋质量}\par 84元按每千克12元算，可买7千克鸡蛋；牛奶36元相当于3千克鸡蛋价钱，扣去即可。\[\begin{aligned} &(100-16)\div12-36\div12\\ ={}&84\div12-3\\ ={}&7-3\\ ={}&4\text{(千克)}\end{aligned}\]答：这次买了4千克鸡蛋。\par
\end{BookAnswer}

% p096_s02
\BookQuestion{五、4（1）}
4. 为鼓励市民节约用电，电力公司制定了以下收费办法：每户每月用电\(100\)千瓦时以内（含\(100\)千瓦时），按\(2\)元/千瓦时收费；超过\(100\)千瓦时的部分，按\(3\)元/千瓦时收费。（6分）\par （1）小亮家10月份用电\(105\)千瓦时，应该交电费多少元？\par
\begin{BookAnswer}
\textbf{方法一：分两档求电费}\par 每千瓦时2元乘100千瓦时；超过的5千瓦时，每千瓦时3元。乘法均按单价乘数量书写。\[\begin{aligned} &2\times100+3\times(105-100)\\ ={}&200+3\times5\\ ={}&200+15\\ ={}&215\text{(元)}\end{aligned}\]答：10月份应交电费215元。\par \textbf{方法二：全部先按高档价，再减多算的钱}\par 先按每千瓦时3元算105千瓦时，前100千瓦时每千瓦时多算1元，扣回。\[\begin{aligned} &3\times105-(3-2)\times100\\ ={}&315-1\times100\\ ={}&315-100\\ ={}&215\text{(元)}\end{aligned}\]答：10月份应交电费215元。\par
\end{BookAnswer}

% p096_s03
\BookQuestion{五、4（2）}
（2）小红家11月份交电费\(260\)元，算一算，她家11月份用电多少千瓦时？\par
\begin{BookAnswer}
\textbf{方法一：先扣前100千瓦时的电费}\par 前100千瓦时收费\(2\times100=200\)元。\(260>200\)，可确定确实用到了每千瓦时3元的高档电量；余60元除以高档单价，求得超过100千瓦时的部分。\[\begin{aligned} &(260-2\times100)\div3+100\\ ={}&(260-200)\div3+100\\ ={}&60\div3+100\\ ={}&20+100\\ ={}&120\text{(千瓦时)}\end{aligned}\]答：11月份用了120千瓦时电。\par \textbf{方法二：把已付费用补成全按3元计费}\par 前100千瓦时每千瓦时少收1元，补上100元后除以每千瓦时3元。\[\begin{aligned} &[260+(3-2)\times100]\div3\\ ={}&[260+1\times100]\div3\\ ={}&[260+100]\div3\\ ={}&360\div3\\ ={}&120\text{(千瓦时)}\end{aligned}\]答：11月份用了120千瓦时电。\par
\end{BookAnswer}

% p096_s04
\BookQuestion{五、5（1）}
5. 某班四年级一班男、女生立定跳远成绩如下。（8分）\par （1）填写下表。\par
\par\smallskip\Needspace{4\baselineskip}\begingroup\small\setlength{\tabcolsep}{3pt}
\noindent\textbf{男生立定跳远成绩（单位：cm）}\par\smallskip
\noindent\begin{tabularx}{\linewidth}{|*{2}{>{\centering\arraybackslash}X|}}\hline
序号 & 成绩\\\hline
1 & 157\\\hline
2 & 159\\\hline
3 & 150\\\hline
4 & 148\\\hline
5 & 128\\\hline
6 & 157\\\hline
7 & 148\\\hline
8 & 125\\\hline
9 & 146\\\hline
10 & 102\\\hline
11 & 149\\\hline
12 & 154\\\hline
13 & 116\\\hline
14 & 138\\\hline
15 & 161\\\hline
16 & 154\\\hline
17 & 149\\\hline
18 & 155\\\hline
19 & 141\\\hline
20 & 143\\\hline
21 & 122\\\hline
22 & 156\\\hline
23 & 154\\\hline
24 & 168\\\hline
\end{tabularx}\par
\endgroup\smallskip
\par\smallskip\Needspace{4\baselineskip}\begingroup\small\setlength{\tabcolsep}{3pt}
\noindent\textbf{女生立定跳远成绩（单位：cm）}\par\smallskip
\noindent\begin{tabularx}{\linewidth}{|*{2}{>{\centering\arraybackslash}X|}}\hline
序号 & 成绩\\\hline
1 & 118\\\hline
2 & 120\\\hline
3 & 145\\\hline
4 & 136\\\hline
5 & 117\\\hline
6 & 148\\\hline
7 & 136\\\hline
8 & 115\\\hline
9 & 137\\\hline
10 & 100\\\hline
11 & 139\\\hline
12 & 144\\\hline
13 & 106\\\hline
14 & 128\\\hline
15 & 161\\\hline
16 & 145\\\hline
17 & 162\\\hline
18 & 148\\\hline
19 & 131\\\hline
20 & 126\\\hline
21 & 102\\\hline
22 & 153\\\hline
23 & 152\\\hline
24 & 142\\\hline
\end{tabularx}\par
\endgroup\smallskip
\par\smallskip\Needspace{4\baselineskip}\begingroup\small\setlength{\tabcolsep}{3pt}
\noindent\textbf{男生立定跳远成绩分段统计表}\par\smallskip
\noindent\begin{tabularx}{\linewidth}{|*{5}{>{\centering\arraybackslash}X|}}\hline
成绩 & 103cm及以下 & 104～127cm & 128～149cm & 150cm及以上\\\hline
人数 &  &  &  & \\\hline
\end{tabularx}\par
\endgroup\smallskip
\par\smallskip\Needspace{4\baselineskip}\begingroup\small\setlength{\tabcolsep}{3pt}
\noindent\textbf{女生立定跳远成绩分段统计表}\par\smallskip
\noindent\begin{tabularx}{\linewidth}{|*{5}{>{\centering\arraybackslash}X|}}\hline
成绩 & 103cm及以下 & 104～127cm & 128～149cm & 150cm及以上\\\hline
人数 &  &  &  & \\\hline
\end{tabularx}\par
\endgroup\smallskip
\begin{BookAnswer}
每名学生的成绩只计入一个分段。核对原表，女生第\(24\)号的成绩是\(142\)厘米，计入\(128\sim149\)厘米这一段。
\[
\begin{array}{c|rrrr|r}
\text{成绩(厘米)}&103\text{及以下}&104\sim127&128\sim149&150\text{及以上}&\text{合计}\\\hline
\text{男生}&1&3&9&11&24\\
\text{女生}&2&6&12&4&24
\end{array}
\]
女生各档学生的序号如下。
\begin{itemize}
\item \(103\)厘米及以下，\allowbreak{}\(10,\ 21\)，共\(2\)人。
\item \(104\sim127\)厘米，\allowbreak{}\(1,\ 2,\ 5,\ 8,\ 13,\ 20\)，共\(6\)人。
\item \(128\sim149\)厘米，\allowbreak{}\(3,\ 4,\ 6,\ 7,\ 9,\ 11,\ 12,\ 14,\ 16,\ 18,\ 19,\ 24\)，共\(12\)人。
\item \(150\)厘米及以上，\allowbreak{}\(15,\ 17,\ 22,\ 23\)，共\(4\)人。
\end{itemize}
核对：男生\(1+3+9+11=24\)（人），女生\(2+6+12+4=24\)（人），全班\(24+24=48\)（人）。
\end{BookAnswer}

% p096_s05
\BookQuestion{五、5（2）}
（2）根据表中数据填空。\par 男生立定跳远在（\hspace{1.8em}）范围的人数最多；女生立定跳远在（\hspace{1.8em}）范围的人数最多；男生和女生，（\hspace{1.8em}）的立定跳远成绩更好一些。\par
\begin{BookAnswer}
男生人数最多的是\(150\)厘米及以上档，共\(11\)人。女生人数最多的是\(128\sim149\)厘米档，共\(12\)人。男生整体成绩更好一些：达到\(150\)厘米及以上的男生有\(11\)人，女生有\(4\)人。达到\(128\)厘米及以上的男生有\(20\)人，女生有\(16\)人。
\end{BookAnswer}

\Needspace{12\baselineskip}

\section{期末综合素养评价（二）}

\Needspace{8\baselineskip}

\pdfbookmark[2]{原书第97页}{orig-97}

\subsection*{原书第97页}

% p097_s01
\BookQuestion{一、1}
期末综合素养评价（二）。时间：90分钟；满分：100分（知识：98分，卷面：2分）。\par 一、填空。（每空1分，共19分）\par 1. 2025年山东省的粮食播种面积是\(\underline{126339000}\)亩，横线上的数改写成用“亿”作单位的数是（\hspace{1.8em}）亩，保留一位小数大约是（\hspace{1.8em}）亿亩。\par
\begin{BookAnswer}
\[126339000\text{亩}=1.26339\text{亿亩}\]原题第一个空后只印“亩”，因此第一个空须填\(1.26339\text{亿}\)。保留一位小数看百分位6，向十分位进1，约为\(1.3\text{亿亩}\)；第二个空后已印“亿亩”，只填\(1.3\)。不能将前一个数写成1.26339亩，也不能在后一个空重复填“亿”。
\end{BookAnswer}

% p097_s02
\BookQuestion{一、2}
2. \(3.925\)由（\hspace{1.8em}）个\(1\)、（\hspace{1.8em}）个\(0.1\)、（\hspace{1.8em}）个\(0.01\)和（\hspace{1.8em}）个\(0.001\)组成，请用▲在数线上表示出这个数。\par
\BookFigure{figures/p097_s02-01.tex}
\begin{BookAnswer}
\(3.925=3+0.9+0.02+0.005\)，依次有3个1、9个0.1、2个0.01、5个0.001。数线上\(3.925-3.9=0.025\)，位于3.9右侧0.025处；若每小格0.005，应向右数5小格。左图红色三角形的顶点对应这个位置。
\BookFigure{figures/p097_s02-02.tex}
\end{BookAnswer}

% p097_s03
\BookQuestion{一、3}
3. 小军在做小数加法时，把其中一个加数\(3.5\)的小数点忽略了，算出来的结果比原来多（\hspace{1.8em}）。\par
\begin{BookAnswer}
忽略3.5的小数点，就把3.5当成35，多加了\(35-3.5=31.5\)。另一个加数没变，所以和也比原来多31.5。
\end{BookAnswer}

% p097_s04
\BookQuestion{一、4}
4. \(32\times(\bigcirc+\triangle)=32\times\underline{\quad}+32\times\underline{\quad}\)，这是运用了\(\underline{\qquad}\)律。\par
\begin{BookAnswer}
把32乘一个和，拆成32分别乘和中的两个加数，再将积相加：\[32\times(\bigcirc+\triangle)=32\times\bigcirc+32\times\triangle.\]依次填圆、三角形、“乘法分配”。
\end{BookAnswer}

% p097_s05
\BookQuestion{一、5}
5. 下面各组直线中，相交的有（\hspace{1.8em}），互相垂直的有（\hspace{1.8em}），互相平行的有（\hspace{1.8em}）。\par
\BookFigure{figures/p097_s05-01.tex}
\begin{BookAnswer}
直线要看无限延长后的关系。①、③、⑤延长后相交；③、⑤的交角是直角，互相垂直；②、④两线方向相同且不相交，互相平行。依次填①③⑤、③⑤、②④。
\end{BookAnswer}

% p097_s06
\BookQuestion{一、6}
6. 如图，将一张长\(10\,\mathrm{cm}\)、宽\(4\,\mathrm{cm}\)的长方形和一张三角形纸交叉摆放，重叠部位是一个梯形，它的高是（\hspace{1.8em}）厘米。如果\(\angle1=102^\circ\)，那么\(\angle2=(\quad)^\circ\)。\par
\BookFigure{figures/p097_s06-01.tex}
\begin{BookAnswer}
重叠梯形的两条底分别在长方形的两条长边上。它们之间的垂直距离就是长方形的宽，所以梯形的高是\(4\)厘米。\par \(\angle1\)与\(\angle2\)组成平角，\allowbreak{}\(\angle2=180^\circ-102^\circ=78^\circ\)。依次填\(4\)、\(78\)。
\end{BookAnswer}

% p097_s07
\BookQuestion{一、7（分割图，问题续下一页）}
7. 学校计划在一块周长\(10\,\mathrm m\)的等腰梯形空地上修建花坛（如右图）。现要将这块梯形空地划分成一个平行四边形和一个三角形的花坛，请你在图中划分一下。若要给三角形花坛四周安装护栏，一共需要护栏（\hspace{1.8em}）米。（题干续下一页）\par
\BookFigure{figures/p097_s07-01.tex}
\begin{BookAnswer}
把梯形顶点顺次记为A、B、C、D，AB为上底，DC为下底。由上底右端B向下底作线段BE，使\(BE\parallel AD\)，E落在DC上，则ABED为平行四边形，BCE为三角形。分割线必须与左腰平行；任意斜着连接不能保证得到平行四边形。左图已校准为这一平行关系。三角形护栏长度的计算见本书原书第98页对应的“一、7（续）”。
\BookFigure{figures/p097_s07-02.tex}
\end{BookAnswer}

\Needspace{8\baselineskip}

\pdfbookmark[2]{原书第98页}{orig-98}

\subsection*{原书第98页}

% p098_s01
\BookQuestion{一、7（续）}
一、7（续）。若要给三角形花坛四周安装护栏，一共需要护栏（\hspace{1.8em}）米。\par
\BookFigure{figures/p098_s01-01.tex}
\begin{BookAnswer}
等腰梯形的两腰相等：\[\begin{aligned} &(10-(2+4))\div2\\ ={}&(10-6)\div2\\ ={}&4\div2\\ ={}&2\text{(米)}\end{aligned}\]平行四边形上下底相等，\allowbreak{}\(DE=AB=2\)米，故\(EC=4-2=2\)米；另有\(BE=AD=2\)米、\(BC=2\)米。三角形三边都是2米：\[\begin{aligned} &2\times3\\ ={}&6\text{(米)}\end{aligned}\]填6。
\end{BookAnswer}

% p098_s02
\BookQuestion{一、8}
8. 四年级一班分成\(8\)个小组进行乒乓球比赛。若采用循环赛制，每组都要和其他所有组赛一场，一共要赛（\hspace{1.8em}）场。若采用淘汰赛制，输一场就淘汰，要决出最后的冠军组，一共要赛（\hspace{1.8em}）场。\par
\BookFigure{figures/p098_s02-01.tex}
\begin{BookAnswer}
循环赛时每组与另外7组各赛一次，\allowbreak{}\(7\times8\)把每场重复计了两次，故列式后先算8除以2，再乘7，运用整数乘除运算的性质。\[\begin{aligned}&7\times8\div2\\={}&7\times(8\div2)\\={}&7\times4\\={}&28\text{(场)}\end{aligned}\]淘汰赛每场淘汰1组，留下冠军须淘汰\(8-1=7\)组，所以共7场。依次填28、7。
\end{BookAnswer}

% p098_s03
\BookQuestion{二、1}
二、选择（将正确答案的序号填在括号里）。（每题2分，共12分）\par 1. 近似数是\(1.45\)的三位小数有（\hspace{1.8em}）个。\par
\noindent A. 5\par
\noindent B. 10\par
\noindent C. 9\par
\noindent D. 4\par
\begin{BookAnswer}
选 C。按四舍五入保留两位小数，\allowbreak{}\(1.445\)到\(1.454\)写成\(1.45\)。本题中\(1.450=1.45\)，按相等处理。\par
小于\(1.45\)的有\(1.445,\ 1.446,\ 1.447,\ 1.448,\ 1.449\)，共\(5\)个。\par
大于\(1.45\)的有\(1.451,\ 1.452,\ 1.453,\ 1.454\)，共\(4\)个。
\[
5+4=9\text{(个)}
\]
\end{BookAnswer}

% p098_s04
\BookQuestion{二、2}
2. 学校为劳动实践园安装防护栏，原来的护栏总长\(20\)米。第一次施工剪掉了\(3.8\)米用于调整转角，第二次又接上了\(4.15\)米延长护栏。现在的护栏和原来相比，（\hspace{1.8em}）。\par
\noindent A. 增加了0.35米\par
\noindent B. 减少了0.35米\par
\noindent C. 增加了7.95米\par
\noindent D. 减少了7.95米\par
\begin{BookAnswer}
选A。第一次剪去3.8米，第二次接上4.15米，净增加\[\begin{aligned} &4.15-3.8\\ ={}&0.35\text{(米)}\end{aligned}\]原来的20米不影响“与原来相比”的增加量。
\end{BookAnswer}

% p098_s05
\BookQuestion{二、3}
3. 下图是三位小朋友在知道了三角形内角和与四边形内角和之后，再探究“六边形内角和是多少”的思考图，你认为正确的是（\hspace{1.8em}）。\par
\noindent A. 甲和乙\par
\noindent B. 甲和丙\par
\noindent C. 乙和丙\par
\noindent D. 甲、乙和丙\par
\BookFigure{figures/p098_s05-01.tex}
\begin{BookAnswer}
选D，甲、乙、丙都正确。甲可分成4个三角形，\allowbreak{}\(180^\circ\times4=720^\circ\)；乙可分成两个四边形，\allowbreak{}\(360^\circ\times2=720^\circ\)；丙分成两个三角形和一个四边形，\allowbreak{}\(180^\circ\times2+360^\circ=720^\circ\)。都覆盖整个六边形且不重叠。
\end{BookAnswer}

% p098_s06
\BookQuestion{二、4}
4. 下列图形，哪一个不能由图①经过一次平移或旋转得到？（\hspace{1.8em}）\par
\noindent A. 图A\par
\noindent B. 图B\par
\noindent C. 图C\par
\noindent D. 图D\par
\BookFigure{figures/p098_s06-01.tex}
\begin{BookAnswer}
选B。平移、旋转会保持图形的形状及边、角的相对排列；B图的缺角方向相当于翻转后的镜像，不能只用一次平移或旋转得到。A、C是旋转得到，D是平移得到。
\end{BookAnswer}

% p098_s07
\BookQuestion{二、5}
5. 为落实国家中小学生体质健康监测工作，对各校进行体质抽测。实验小学\(8\)人的平均成绩是\(92\)分，老师核对成绩后，发现有一位同学的得分少了\(32\)分，调整后实验小学的平均成绩会（\hspace{1.8em}）。\par
\noindent A. 减少32分\par
\noindent B. 增加32分\par
\noindent C. 减少4分\par
\noindent D. 增加4分\par
\begin{BookAnswer}
选 D。补记这位同学的\(32\)分后，\allowbreak{}\(8\)人的总分增加\(32\)分。把增加的分数平均分给\(8\)人。
\[
\begin{aligned}
&32\div8\\
={}&4\text{(分)}
\end{aligned}
\]
调整后的平均成绩比原来增加\(4\)分。
\end{BookAnswer}

\Needspace{8\baselineskip}

\pdfbookmark[2]{原书第99页}{orig-99}

\subsection*{原书第99页}

% p099_s01
\BookQuestion{二、6}
6. 一张正方形纸，至少剪（\hspace{1.8em}）次，会得到下面的轴对称图形。\par
\noindent A. 1\par
\noindent B. 2\par
\noindent C. 3\par
\noindent D. 4\par
\BookFigure{figures/p099_s01-01.tex}
\begin{BookAnswer}
选A，至少剪1次。可先将正方形对折，使需要剪去的角叠在一起，再一次剪去重合的小角；展开后形成所示轴对称图形。这里统计的是剪的次数，不是对折次数。
\end{BookAnswer}

% p099_s02
\BookQuestion{三、1}
三、计算。（共17分）\par 1. 直接写得数。（8分）\par \(4.3+5.6=\quad9.74-3.4=\quad0.58+0.3=\quad8.25-2.25=\)\par \(2.4\times5=\quad1.25\times8=\quad45+240\div8=\quad12\times5-30=\)\par
\begin{BookAnswer}
按原排列，第一行\(9.9,6.34,0.88,6\)；第二行\(12,10,75,30\)。\par \(45+240\div8\)先除后加，\allowbreak{}\(240\div8=30\)，\allowbreak{}\(45+30=75\)；\allowbreak{}\(12\times5-30=60-30=30\)。其余小数加减要对齐小数点。\par 本阶段对\(2.4\times5\)可按计数单位理解：24个十分之一的5倍是120个十分之一，即12；\allowbreak{}\(1.25\times8\)是125个百分之一的8倍，等于1000个百分之一，即10。不需套用尚未系统学习的一般小数乘法竖式。
\end{BookAnswer}

% p099_s03
\BookQuestion{三、2（1）}
2. 计算下列各题，能简便计算的要简便计算。（9分）\par \(45.2-13.64-6.36\)\par
\begin{BookAnswer}
将连续减去的两数先合并。\[\begin{aligned} &45.2-13.64-6.36\\ ={}&45.2-(13.64+6.36)\\ ={}&45.2-20\\ ={}&25.2\end{aligned}\]
\end{BookAnswer}

% p099_s04
\BookQuestion{三、2（2）}
\(27\times99-27\times49\)\par
\begin{BookAnswer}
反用乘法分配律。\[\begin{aligned} &27\times99-27\times49\\ ={}&27\times(99-49)\\ ={}&27\times50\\ ={}&1350\end{aligned}\]
\end{BookAnswer}

% p099_s05
\BookQuestion{三、2（3）}
\(960\div[(32+16)\div4]\)\par
\begin{BookAnswer}
除以一个商，可以先除以被除数，再乘除数；先算\(960\div48\)较简便。\[\begin{aligned}&960\div[(32+16)\div4]\\={}&960\div[48\div4]\\={}&960\div48\times4\\={}&20\times4\\={}&80\end{aligned}\]
\end{BookAnswer}

% p099_s06
\BookQuestion{四、1}
四、操作。（共15分）\par 1. 为了打造更舒适的校园环境，学校重新进行了校园规划，打算从主干道修建一条最短的小路直通图书馆，在主干道的北面过点\(A\)修一条东西方向的健身步道，请你根据要求完善规划图。（6分）\par
\BookFigure{figures/p099_s06-01.tex}
\begin{BookAnswer}
从图书馆点向主干道作垂线段，垂足落在主干道上并标直角；因为点到直线的垂线段最短，这就是最短小路。\par 北箭头朝上，东西方向为水平方向；过A作水平直线，且A点要在线上，就是所求健身步道。
\BookFigure{figures/p099_s06-02.tex}
\end{BookAnswer}

% p099_s07
\BookQuestion{四、2（1）}
2. 画一画。（9分）\par （1）画出图形①的另一半，使它成为一个轴对称图形。\par
\BookFigure{figures/p099_s07-01.tex}
\begin{BookAnswer}
以图中已给出的竖直虚线为对称轴。按每个顶点到轴的水平格数，在轴的另一侧找到对应点，再按图形轮廓依次连接。轴上的点位置不变，保留原有的虚线。
\BookFigure{figures/p099_s07-02.tex}
\end{BookAnswer}

% p099_s08
\BookQuestion{四、2（2）}
（2）将图形②先绕\(O\)点顺时针旋转\(90^\circ\)，再将旋转后的图形向右平移\(3\)格，分别画出旋转后和平移后的图形。\par
\BookFigure{figures/p099_s08-01.tex}
\begin{BookAnswer}
以\(O\)为旋转中心，将另两个顶点分别绕\(O\)顺时针旋转\(90^\circ\)，与\(O\)连成第一个红色三角形。再把该三角形的三个顶点都向右移动\(3\)格，依次连接，得到第二个红色三角形。
\BookFigure{figures/p099_s08-02.tex}
\end{BookAnswer}

\Needspace{8\baselineskip}

\pdfbookmark[2]{原书第100页}{orig-100}

\subsection*{原书第100页}

% p100_s01
\BookQuestion{五、1（1）}
五、解决问题。（共35分）\par 1. 社区开展助农直播活动，助力农户售卖水果与干货。\par （1）农户上午卖出\(15\)箱苹果，每箱\(32\)元；卖出\(12\)箱梨，每箱\(24\)元。农户上午一共收入多少元？（5分）\par
\begin{BookAnswer}
\textbf{方法一：分别算两种水果收入}\par 每箱32元乘15箱；每箱24元乘12箱，再相加。\[\begin{aligned} &32\times15+24\times12\\ ={}&480+288\\ ={}&768\text{(元)}\end{aligned}\]答：上午共收入768元。\par \textbf{方法二：全部先按每箱24元，再补差价}\par 共27箱先按每箱24元算，15箱苹果每箱还应补8元。\[\begin{aligned} &24\times(15+12)+(32-24)\times15\\ ={}&24\times27+8\times15\\ ={}&648+120\\ ={}&768\text{(元)}\end{aligned}\]答：上午共收入768元。\par
\end{BookAnswer}

% p100_s02
\BookQuestion{五、1（2）}
（2）农户下午卖出\(80\)袋核桃，每袋售价\(45\)元，每袋核桃需缴纳平台服务费\(2\)元。农户全天实际一共收入多少元？（5分）\par
\begin{BookAnswer}
上午实际收入是\(768\)元。下午每卖一袋核桃，要从售价中扣除\(2\)元平台服务费。\par
\textbf{方法一：先求每袋实际收入}\par
每袋实际收入是\((45-2)\)元，卖出\(80\)袋，再加上上午的\(768\)元。
\[
\begin{aligned}
&768+(45-2)\times80\\
={}&768+43\times80\\
={}&768+3440\\
={}&4208\text{(元)}
\end{aligned}
\]
答：农户全天实际一共收入\(4208\)元。\par
\textbf{方法二：下午销售额扣除平台服务费}\par
下午销售额为\(45\times80\)元，总服务费为\(2\times80\)元。两项都有\(80\)份，可以用乘法分配律合并。
\[
\begin{aligned}
&768+45\times80-2\times80\\
={}&768+(45-2)\times80\\
={}&768+43\times80\\
={}&768+3440\\
={}&4208\text{(元)}
\end{aligned}
\]
答：农户全天实际一共收入\(4208\)元。
\end{BookAnswer}

% p100_s03
\BookQuestion{五、2}
2. 文具店售卖一款可换芯中性笔，笔杆长\(0.14\)米，配套的笔帽长\(0.05\)米。将笔帽套在笔杆上时，总长度是\(0.15\)米。请你计算一下笔帽和笔杆重叠部分的长度是多少厘米。（7分）\par
\BookFigure{figures/p100_s03-01.tex}
\begin{BookAnswer}
\textbf{方法一：两部分全长之和减套好后的总长}\par 笔杆和笔帽分别量长后相加，重叠部分被多算了一次，减去总长即可。\[\begin{aligned} &0.14+0.05-0.15\\ ={}&0.19-0.15\\ ={}&0.04\text{(米)}\end{aligned}\]\[0.04\text{米}=4\text{厘米}\]答：重叠部分长4厘米。\par \textbf{方法二：从笔帽长中扣掉露出的笔帽长}\par 套好后比笔杆长多出0.01米，这0.01米是笔帽没有重叠的部分。\[\begin{aligned} &0.05-(0.15-0.14)\\ ={}&0.05-0.01\\ ={}&0.04\text{(米)}\end{aligned}\]\[0.04\text{米}=4\text{厘米}\]答：重叠部分长4厘米。\par
\end{BookAnswer}

% p100_s04
\BookQuestion{五、3}
3. 青岛到上海的高速公路全长约\(720\)千米，甲、乙两辆客车分别从青岛和上海同时相对开出，行驶\(3\)小时后，还相距\(240\)千米。甲客车的速度是\(90\)千米/时，求乙客车的速度。（先画图表示数量关系，再解答。）（9分）\par
\begin{BookAnswer}
图中总路程720千米被分成甲车3小时所行路程、两车相距240千米、乙车3小时所行路程三段。\par \textbf{方法一：先求两车速度之和}\par 先求3小时两车一共行多少千米，除以3求速度和，再减甲车速度。\[\begin{aligned} &(720-240)\div3-90\\ ={}&480\div3-90\\ ={}&160-90\\ ={}&70\text{(千米/时)}\end{aligned}\]答：乙车每小时约行70千米。\par \textbf{方法二：先求乙车行驶路程}\par 从720千米中依次扣掉尚余240千米及甲车所行270千米，余下是乙车3小时行驶路程。\[\begin{aligned} &(720-240-90\times3)\div3\\ ={}&(720-240-270)\div3\\ ={}&(480-270)\div3\\ ={}&210\div3\\ ={}&70\text{(千米/时)}\end{aligned}\]答：乙车每小时约行70千米。\par \par 代回线段图检验：甲、乙已行路程与尚余路程相加为\(90\times3+70\times3+240=270+210+240=720\)千米，恰好等于两地总路程。
\BookFigure{figures/p100_s04-01.tex}
\end{BookAnswer}

% p100_s05
\BookQuestion{五、4（1）（2）}
4. 下面是四（1）班同学的综合素质成绩（单位：分）。（9分）\par \(58\;95\;86\;75\;84\;87\;59\;90\;77\;98\;93\;87\;58\;62\;77\;86\)\par \(98\;68\;90\;87\;96\;95\;98\;88\;90\;77\;95\;84\;99\;67\;78\;88\)\par \(90\;91\;86\;93\)\par （1）请根据以上数据，完成下面的统计表。\par 备注：85～100分为优秀；70～84分为良好；60～69分为合格；60分以下为不合格。\par （2）你认为该班级综合素质是怎样的？说明理由。\par
\par\smallskip\Needspace{4\baselineskip}\begingroup\small\setlength{\tabcolsep}{3pt}
\noindent\begin{tabularx}{\linewidth}{|*{6}{>{\centering\arraybackslash}X|}}\hline
等级 & 合计 & 优秀 & 良好 & 合格 & 不合格\\\hline
人数 &  &  &  &  & \\\hline
\end{tabularx}\par
\endgroup\smallskip
\begin{BookAnswer}
按给定等级边界逐人成绩登记：\(85\sim100\)分23人；\allowbreak{}\(70\sim84\)分7人；\allowbreak{}\(60\sim69\)分3人；60分以下3人。\par 合计\(23+7+3+3=36\)人。依次填合计36、优秀23、良好7、合格3、不合格3；边界85归优秀，70和84归良好，60和69归合格。\par \textbf{（2）说明：}整体情况较好：36人中优秀23人，23大于36的一半；优秀与良好共\(23+7=30\)人。也应关注3名不合格同学，查找具体困难并提供帮助。说明必须引用表中数据，不能只写“好”或“不好”。
\end{BookAnswer}

\Needspace{12\baselineskip}

\section{期末综合素养评价（三）}

\Needspace{8\baselineskip}

\pdfbookmark[2]{原书第101页}{orig-101}

\subsection*{原书第101页}

% p101_s01
\BookQuestion{一、1}
期末综合素养评价（三）。时间：90分钟；满分：100分（知识：98分，卷面：2分）。\par 一、填空。（每空1分，共15分）\par 1. 一个数是由\(6\)个一、\(3\)个十分之一和\(5\)个百分之一组成，这个数写作（\hspace{1.8em}），读作（\hspace{1.8em}），精确到十分位是（\hspace{1.8em}）。\par
\begin{BookAnswer}
6个一、3个十分之一、5个百分之一合为\(6+0.3+0.05=6.35\)，读作“六点三五”。精确到十分位看百分位5，向十分位进1得\(6.4\)。
\end{BookAnswer}

% p101_s02
\BookQuestion{一、2}
2. 把\(128+147=275\)，\allowbreak{}\(275\div25=11\)，\allowbreak{}\(240\times11=2640\)这三个算式写成一个综合算式是（\hspace{1.8em}）。\par
\begin{BookAnswer}
先求和\(128+147\)，再除以25，最后乘240。综合算式为\[240\times[(128+147)\div25].\]外层240乘括号，不能把147单独除以25。
\end{BookAnswer}

% p101_s03
\BookQuestion{一、3}
3. 在○里填上“＞”“＜”或“＝”。\par \(0.060\) ○ \(0.06\)；\allowbreak{}\(7\)角\(5\)分 ○ \(0.8\)元；\allowbreak{}\(7.2\)米 ○ \(7\)米\(2\)厘米；\allowbreak{}\(65\times100-1\) ○ \(65\times99\)；\allowbreak{}\(1000-156-244\) ○ \(1000-(156+244)\)\par
\begin{BookAnswer}
（1）\(0.060=0.06\)，小数末尾添0不改变大小，填\(=\)。\par （2）先统一单位：\[7\text{角}5\text{分}=0.75\text{元}\]\(0.75<0.8\)，填\(<\)。\par （3）先统一单位：\[7\text{米}2\text{厘米}=7.02\text{米}\]\(7.2>7.02\)，填\(>\)。\par （4）\(65\times100-1=6499\)，\allowbreak{}\(65\times99=6435\)，填\(>\)。\par （5）两边都是\(1000-400=600\)，填\(=\)。
\end{BookAnswer}

% p101_s04
\BookQuestion{一、4}
4. 一个等腰三角形，它的一个底角是\(45^\circ\)，它的顶角是（\hspace{1.8em}）°，按角分类，这个三角形是（\hspace{1.8em}）三角形。\par
\begin{BookAnswer}
两个底角都是45度，顶角为\[\begin{aligned} &180^\circ-45^\circ\times2\\ ={}&180^\circ-90^\circ\\ ={}&90^\circ\end{aligned}\]有一个直角，属直角三角形；依次填90、直角。
\end{BookAnswer}

% p101_s05
\BookQuestion{一、5}
5. 把一张正方形纸沿相同的方向对折两次，打开后出现两条折痕，折痕的位置关系是（\hspace{1.8em}）。\par
\begin{BookAnswer}
沿相同方向对折两次，两条折痕方向相同、保持间隔，互相平行。若横竖两个方向各折一次，才会出现垂直折痕。
\end{BookAnswer}

% p101_s06
\BookQuestion{一、6}
6. 如图，一张三角形和一张平行四边形纸片。把三角形的各个角撕下来拼在一起，可以组成一个（\hspace{1.8em}）角，把平行四边形的各个角撕下来拼在一起可以组成一个（\hspace{1.8em}）角。\par
\BookFigure{figures/p101_s06-01.tex}
\begin{BookAnswer}
三角形三个角拼在一起的总度数是\(180^\circ\)，可组成平角；平行四边形四个角合为\(360^\circ\)，可组成周角。依次填“平”“周”。
\end{BookAnswer}

% p101_s07
\BookQuestion{一、7}
7. 四年级一班参加书法社团的有\(18\)人，参加绘画社团的有\(16\)人，两个社团都参加的有\(5\)人，只参加一个社团的有（\hspace{1.8em}）人。\par
\begin{BookAnswer}
只参加书法的有\(18-5=13\)人，只参加绘画的有\(16-5=11\)人，共\(13+11=24\)人。\par 也可\(18+16-5\times2=24\)；两项都参加的人在18和16中各被算一次，而本问要求他们一次也不计，因此要减两次5。
\end{BookAnswer}

\Needspace{8\baselineskip}

\pdfbookmark[2]{原书第102页}{orig-102}

\subsection*{原书第102页}

% p102_s01
\BookQuestion{二、1（题干跨页）}
二、选择（将正确答案的序号填在括号里）。（每题2分，共14分）\par 1. 万老师用三根木条制作了一个衣架（如图），这个衣架的三根木条的长度可能是（\hspace{1.8em}）。\par
\noindent A. 2.1dm，2.1dm，[首位不清].2dm\par
\noindent B. 2dm，3dm，3dm\par
\noindent C. 2.4dm，2.4dm，4.2dm\par
\BookFigure{figures/p102_s01-01.tex}
\begin{BookAnswer}
至少B、C两项均能围成三角形，原单选设置不唯一。B两短边之和\(2+3=5>3\)；C两短边之和\(2.4+2.4=4.8>4.2\)，均满足“两短边之和大于最长边”。A的末长首位扫描缺损，不能凭猜读排除；若末长为7.2（或4.2）分米，则\(2.1+2.1\le\)末长，不能围成。B、C均符合的结论不受A模糊影响。
\end{BookAnswer}

% p102_s02
\BookQuestion{二、2}
2. 下列说法正确的是（\hspace{1.8em}）。\par ①\(6.995\)精确到百分位是\(7.00\)。\par ②右图竖式的计算过程中运用了乘法分配律。\par ③\(0.2\)与\(0.4\)之间只有一个小数。\par ④相机支架一般都是三角形状，这是利用了三角形的稳定性。\par
\noindent A. ①②④\par
\noindent B. ①③④\par
\noindent C. ①④\par
\BookFigure{figures/p102_s02-01.tex}
\begin{BookAnswer}
选A，即①②④正确。\par ①\(6.995\)精确到百分位，千分位5进位后是\(7.00\)，两个0要保留。\par ②\(48\times37=48\times(30+7)=48\times30+48\times7\)，竖式两行积体现乘法分配律。\par ③错误，\allowbreak{}\(0.2\)和\(0.4\)之间有无数个小数，例如0.21、0.22，仍可继续细分。\par ④三角形有稳定性，相机支架常用三角形结构。
\end{BookAnswer}

% p102_s03
\BookQuestion{二、3}
3. 优优、佳佳、浩浩三名同学在一次跳绳比赛中的平均成绩是\(125\)个，乐乐跳了\(145\)个，他们四人跳的平均成绩是（\hspace{1.8em}）个。\par
\noindent A. 125\par
\noindent B. 130\par
\noindent C. 145\par
\begin{BookAnswer}
选B，先求三人的跳绳总数，再加乐乐的145个，除以4：\[\begin{aligned} &(125\times3+145)\div4\\ ={}&(375+145)\div4\\ ={}&520\div4\\ ={}&130\text{(个)}\end{aligned}\]\par 也可移多补少：乐乐比原平均数多\(145-125=20\)个，把多出的20个平均分给4人。\[\begin{aligned}&125+(145-125)\div4\\={}&125+20\div4\\={}&125+5\\={}&130\text{(个)}\end{aligned}\]原来的3人不是每人都跳125个，但他们的总数等于\(125\times3\)，所以仍可按这个基准分配。
\end{BookAnswer}

% p102_s04
\BookQuestion{二、4}
4. 下列的关系图中，能正确表示同一平面内两条直线位置关系的是（\hspace{1.8em}）。\par
\noindent A. 垂直属于相交，平行与相交并列\par
\noindent B. 垂直属于平行\par
\noindent C. 相交、平行、垂直并列\par
\BookFigure{figures/p102_s04-01.tex}
\begin{BookAnswer}
选A。同一平面内，两条不同直线的位置关系分成平行与相交；互相垂直是相交中的特殊情况，应包含在“相交”部分内，不能与相交并列为互不包含的类别。
\end{BookAnswer}

% p102_s05
\BookQuestion{二、5}
5. 李华在课中研究三角形，用直尺、圆规绘制了下面的三幅图，你知道他在研究什么问题吗？（\hspace{1.8em}）\par
\noindent A. 三角形的特征\par
\noindent B. 三角形的内角和\par
\noindent C. 三角形三边关系\par
\noindent D. 三角形的周长\par
\BookFigure{figures/p102_s05-01.tex}
\begin{BookAnswer}
选C。三幅图用两点间距离作圆弧，逐步比较或确定第三边是否能与另外两边围成三角形，研究的是三角形的三边关系，不是内角和或周长。
\end{BookAnswer}

% p102_s06
\BookQuestion{二、6（接下一页选项）}
6. 短跑赛场上有两组队员正在准备比赛，采取三局两胜制。如果比赛中每个人都正常发挥，根据“田忌赛马”的策略，按照（\hspace{1.8em}）方案安排出场顺序，二组才能获胜。（选项与成绩见下一页）\par
\begin{BookAnswer}
选C。按丙对A、乙对C、甲对B：丙24秒对A15秒，二组输；乙21秒对C23秒，二组赢；甲16秒对B20秒，二组赢。短跑用时少者胜，二组赢两场，符合三局两胜。其他三个方案都只能赢一场。
\end{BookAnswer}

\Needspace{8\baselineskip}

\pdfbookmark[2]{原书第103页}{orig-103}

\subsection*{原书第103页}

% p103_s01
\BookQuestion{二、6（选项在本页）}
二、6（续）。请结合上一页题干及本页成绩表，选择使二组在三局两胜制中获胜的方案。\par
\noindent A. 丙—A，乙—B，甲—C\par
\noindent B. 丙—B，乙—C，甲—A\par
\noindent C. 丙—A，乙—C，甲—B\par
\noindent D. 丙—C，乙—A，甲—B\par
\par\smallskip\Needspace{4\baselineskip}\begingroup\small\setlength{\tabcolsep}{3pt}
\noindent\begin{tabularx}{\linewidth}{|*{4}{>{\centering\arraybackslash}X|}}\hline
 & 第一名 & 第二名 & 第三名\\\hline
一组 & A：15秒 & B：20秒 & C：23秒\\\hline
二组 & 甲：16秒 & 乙：21秒 & 丙：24秒\\\hline
\end{tabularx}\par
\endgroup\smallskip
\begin{BookAnswer}
选C，丙对A、乙对C、甲对B，二组依次输、赢、赢。具体用时比较为\(24>15\)、\(21<23\)、\(16<20\)，所以二组赢两场。
\end{BookAnswer}

% p103_s02
\BookQuestion{二、7}
7. 小军在探索四边形的内角和时，采用了如下图所示的分割转化法，下列选项中，（\hspace{1.8em}）体现了他的推算过程。\par
\noindent A. \(4\times180^\circ-360^\circ\)\par
\noindent B. \(\angle1+\angle2+\angle3+\angle4\)\par
\noindent C. \(3\times180^\circ-180^\circ\)\par
\BookFigure{figures/p103_s02-01.tex}
\begin{BookAnswer}
选A。连接内点E与四个顶点，得到4个三角形；四个三角形内角和共\(180^\circ\times4\)，其中E周围的周角\(360^\circ\)不是四边形内角，须减去。\[\begin{aligned}&180^\circ\times4-360^\circ\\={}&720^\circ-360^\circ\\={}&360^\circ\end{aligned}\]B只合计了图中部分小角，漏掉了其他内角部分。
\end{BookAnswer}

% p103_s03
\BookQuestion{三、1}
三、计算。（共30分）\par 1. 直接写得数。（8分）\par \(3.26+0.4=\quad4.37+2.7=\quad5.6-4=\quad62+38\times0=\)\par \(0.65+0.35=\quad0.6\div100=\quad8-2.68=\quad4.32\times1000=\)\par
\begin{BookAnswer}
按原题行列依次为：\par 第一行\(3.66,7.07,1.6,62\)；\par 第二行\(1,0.006,5.32,4320\)。\par 要点：\(62+38\times0=62+0=62\)，先乘后加；\allowbreak{}\(0.6\div100=0.006\)；\allowbreak{}\(4.32\times1000=4320\)。小数加减对齐小数点。
\end{BookAnswer}

% p103_s04
\BookQuestion{三、2（1）}
2. 列竖式计算。（带★的要验算）（13分）\par \(9.27+15.9=\textcolor{answerred}{25.17}\)\par
\begin{BookAnswer}
把15.9末尾补0，与9.27对齐小数点。\[\begin{array}{r}9.27\\+15.90\\\hline25.17\end{array}\]百分位\(7+0=7\)，十分位\(2+9=11\)，满十进一，结果25.17。
\end{BookAnswer}

% p103_s05
\BookQuestion{三、2（2）}
\(25.7+16.4=\textcolor{answerred}{42.1}\)\par
\begin{BookAnswer}
小数点对齐，十分位\(7+4=11\)，向个位进1：\[\begin{array}{r}25.7\\+16.4\\\hline42.1\end{array}\]
\end{BookAnswer}

% p103_s06
\BookQuestion{三、2（3）}
\(10-4.03=\textcolor{answerred}{5.97}\)\par
\begin{BookAnswer}
把10写作10.00，以便按位退位相减：\[\begin{array}{r}10.00\\-4.03\\\hline5.97\end{array}\]检验可算\(5.97+4.03=10.00\)。
\end{BookAnswer}

% p103_s07
\BookQuestion{三、2（4）}
★\(18.2-6.78=\textcolor{answerred}{11.42}\)（要验算）\par
\begin{BookAnswer}
把18.2写作18.20，小数点对齐：\[\begin{array}{r}18.20\\-6.78\\\hline11.42\end{array}\]\par 验算用差加减数，恢复被减数：\[\begin{array}{r}11.42\\+6.78\\\hline18.20\end{array}\]
\end{BookAnswer}

% p103_s08
\BookQuestion{三、3（1）}
3. 计算下列各题，能简便计算的要简便计算。（9分）\par \(45.01+6.1+0.99+2.9\)\par
\begin{BookAnswer}
凑整相加。\[\begin{aligned} &45.01+6.1+0.99+2.9\\ ={}&(45.01+0.99)+(6.1+2.9)\\ ={}&46+9\\ ={}&55\end{aligned}\]
\end{BookAnswer}

% p103_s09
\BookQuestion{三、3（2）}
\(45\times98\)\par
\begin{BookAnswer}
简便方法：98看作100少2。\[\begin{aligned} &45\times98\\ ={}&45\times(100-2)\\ ={}&45\times100-45\times2\\ ={}&4500-90\\ ={}&4410\end{aligned}\]同一分配律的另一种拆法（不另计解法）：45拆成50少5。\[\begin{aligned} &45\times98\\ ={}&(50-5)\times98\\ ={}&50\times98-5\times98\\ ={}&4900-490\\ ={}&4410\end{aligned}\]
\end{BookAnswer}

% p103_s10
\BookQuestion{三、3（3）}
\(101\times101-101\)\par
\begin{BookAnswer}
提取101，减去一份101。\[\begin{aligned} &101\times101-101\\ ={}&101\times(101-1)\\ ={}&101\times100\\ ={}&10100\end{aligned}\]
\end{BookAnswer}

% p103_s11
\BookQuestion{四、（1）}
四、操作。（共12分）\par （1）画出三角形\(ABC\)中\(AC\)边上的高。\par
\BookFigure{figures/p103_s11-01.tex}
\begin{BookAnswer}
AC作为底，应从与AC相对的顶点B作高；过B作\(AC\)的垂线，与AC交于D。高是线段BD，不是“直线BD”，D处要有直角标记。左图D已落在AC上。
\BookFigure{figures/p103_s11-02.tex}
\end{BookAnswer}

% p103_s12
\BookQuestion{四、（2）}
（2）以\(AB\)边所在直线为对称轴，画出轴对称图形的另一半。\par
\BookFigure{figures/p103_s12-01.tex}
\begin{BookAnswer}
以AB所在直线为对称轴，A、B在轴上保持不动。C的对应点\(C_1\)位于轴的右侧，到AB的格数等于C到AB的格数；连\(AC_1,BC_1\)补全图形；\allowbreak{}\(CC_1\)与\(AB\)垂直。
\BookFigure{figures/p103_s12-02.tex}
\end{BookAnswer}

% p103_s13
\BookQuestion{四、（3）}
（3）把三角形\(ABC\)先绕点\(A\)逆时针旋转\(90^\circ\)，再向右平移\(4\)格。\par
\BookFigure{figures/p103_s13-01.tex}
\begin{BookAnswer}
先绕A逆时针旋转\(90^\circ\)：A不动，其余两个顶点到A的距离不变，方向逆时针转直角；连成旋转后的三角形。再把所得三个顶点都向右移动4格，依次连接得到终图。平移的4格应从“旋转后的各顶点”起数，不能仍从原图顶点起数。
\BookFigure{figures/p103_s13-02.tex}
\end{BookAnswer}

\Needspace{8\baselineskip}

\pdfbookmark[2]{原书第104页}{orig-104}

\subsection*{原书第104页}

% p104_s01
\BookQuestion{五、1}
五、解决问题。（共27分）\par 1. 小红的储蓄罐里原来有\(12.3\)元，昨天买笔用去\(7.5\)元，今天妈妈又给了她\(2.3\)元，现在小红的储蓄罐里有多少钱？（5分）\par
\begin{BookAnswer}
\textbf{方法一：按事情发生顺序计算}\par 先从原有12.3元扣昨天买笔的7.5元，再加妈妈今天给的2.3元。\[\begin{aligned} &12.3-7.5+2.3\\ ={}&4.8+2.3\\ ={}&7.1\text{(元)}\end{aligned}\]答：现在储蓄罐里有7.1元。\par \textbf{方法二：先合计拥有过的钱，再扣支出}\par 把原有和新收入先相加，最后减买笔支出。\[\begin{aligned} &12.3+2.3-7.5\\ ={}&14.6-7.5\\ ={}&7.1\text{(元)}\end{aligned}\]答：现在储蓄罐里有7.1元。\par
\end{BookAnswer}

% p104_s02
\BookQuestion{五、2}
2. 利群商场举办“迎十一”促销活动，下图的袜子每双\(4.68\)元，王阿姨买了\(12\)双，花了多少钱？（5分）\par 促销：买五送一。\par
\BookFigure{figures/p104_s02-01.tex}
\begin{BookAnswer}
\textbf{方法一：按六双一组，求付费双数}\par 买5双送1双，6双为一组；12双有2组，每组付5双的钱。\[\begin{aligned}&4.68\times\{5\times[12\div(5+1)]\}\\={}&4.68\times\{5\times[12\div6]\}\\={}&4.68\times\{5\times2\}\\={}&4.68\times10\\={}&46.8\text{(元)}\end{aligned}\]答：12双袜子共花46.8元。\par \textbf{方法二：化成整数分，再扣赠品标价}\par \[4.68\text{元}=468\text{分}\]12双中有2双赠送，列总标价减赠品标价，再用乘法分配律合并为10份，避免先算两个大乘积。\[\begin{aligned}&468\times12-468\times[12\div(5+1)]\\={}&468\times12-468\times[12\div6]\\={}&468\times12-468\times2\\={}&468\times(12-2)\\={}&468\times10\\={}&4680\text{(分)}\end{aligned}\]\[4680\text{分}=46.8\text{元}\]答：12双袜子共花46.8元。
\end{BookAnswer}

% p104_s03
\BookQuestion{五、3}
3. 商店每个足球售价\(108\)元，每个篮球售价\(92\)元。学校需要购买足球和篮球各\(15\)个，一共需要花多少钱？（5分）\par
\begin{BookAnswer}
足球和篮球各15个，可把各一个配成一套，先将每套单价108与92凑200。分别列两类球总价时，利用乘法分配律合并共同的15，也化为同一简便式。\par \[\begin{aligned}&108\times15+92\times15\\={}&(108+92)\times15\\={}&200\times15\\={}&3000\text{(元)}\end{aligned}\]答：一共需要3000元。
\end{BookAnswer}

% p104_s04
\BookQuestion{五、4}
4. 同学们周末参加户外徒步活动。上山时走\(1080\)米用了\(18\)分钟，按原路返回时只用了\(15\)分钟。上山的速度比下山的速度每分钟慢多少？（6分）\par
\begin{BookAnswer}
\textbf{方法一：分别求下山和上山速度}\par 路程相同，速度等于路程除以时间，再求速度差。\[\begin{aligned} &1080\div15-1080\div18\\ ={}&72-60\\ ={}&12\text{(米/分)}\end{aligned}\]答：上山速度比下山速度每分钟慢12米。\par \textbf{方法二：比较同样15分钟所行路程}\par 上山速度为1080除以18等于60米/分；按这个速度15分钟只能走900米，比下山路程少180米，再平均到15分钟。\[\begin{aligned} &(1080-60\times15)\div15\\ ={}&(1080-900)\div15\\ ={}&180\div15\\ ={}&12\text{(米/分)}\end{aligned}\]答：上山速度比下山速度每分钟慢12米。\par
\end{BookAnswer}

% p104_s05
\BookQuestion{五、5}
5. 甲、乙两列车同时从两地出发，相向而行。\(2\)小时后两车第一次相距\(200\)千米。甲车每小时行\(220\)千米，乙车每小时行\(280\)千米，两地之间的距离是多少千米？（6分）\par
\begin{BookAnswer}
第一次相距200千米表示尚未相遇，已行路程和要加尚余200千米。可先求两车速度和；分别列两车路程时，利用乘法分配律提取共同的2，将220与280凑500。\par \[\begin{aligned}&220\times2+280\times2+200\\={}&(220+280)\times2+200\\={}&500\times2+200\\={}&1000+200\\={}&1200\text{(千米)}\end{aligned}\]答：两地之间相距1200千米。
\end{BookAnswer}

\Needspace{12\baselineskip}

\section{期末综合素养评价（四）}

\Needspace{8\baselineskip}

\pdfbookmark[2]{原书第105页}{orig-105}

\subsection*{原书第105页}

% p105_s01
\BookQuestion{一、1}
期末综合素养评价（四）。时间：90分钟；满分：100分（知识：98分，卷面：2分）。\par 一、填空。（每空1分，共21分）\par 1. \(9.08\)读作（\hspace{1.8em}），它的计数单位是（\hspace{1.8em}），它有（\hspace{1.8em}）个这样的计数单位，再添上（\hspace{1.8em}）个这样的计数单位就是最小的两位数。\par
\begin{BookAnswer}
\(9.08\)读作“九点零八”；最低位为百分位，计数单位为\(0.01\)。9个一等于900个百分之一，再加8个百分之一，共\(900+8=908\)个百分之一。最小两位数10含1000个百分之一，还须添\(1000-908=92\)个。依次填九点零八、0.01、908、92。
\end{BookAnswer}

% p105_s02
\BookQuestion{一、2}
2. \(6\)千克\(50\)克＝（\hspace{1.8em}）克；\allowbreak{}\(15.6\)元＝（\hspace{1.8em}）元（\hspace{1.8em}）角；\allowbreak{}\(900\)克＝（\hspace{1.8em}）千克；\allowbreak{}\(3.07\)千米＝（\hspace{1.8em}）千米（\hspace{1.8em}）米。\par
\begin{BookAnswer}
按单位换算：\[6\text{千克}50\text{克}=(6000+50)\text{克}=6050\text{克}\]\[15.6\text{元}=15\text{元}6\text{角}\]\[900\text{克}=0.9\text{千克}\]\[3.07\text{千米}=3\text{千米}70\text{米}\]填6050、15、6、0.9、3、70，原空后的单位不重复填写。
\end{BookAnswer}

% p105_s03
\BookQuestion{一、3}
3. 用\(1,0,2,7\)和小数点组成一个三位小数，这个小数最大是（\hspace{1.8em}），最小是（\hspace{1.8em}）。\par
\begin{BookAnswer}
用1、0、2、7各一次组成三位小数，整数部分取最大数字7，余下数字从大到小排为210，最大\(7.210\)。最小时整数部分取0，小数部分从小到大排为127，最小\(0.127\)。最大值末尾的0仍占千分位。
\end{BookAnswer}

% p105_s04
\BookQuestion{一、4}
4. 慢跑\(10\)千米大约消耗\(631\)千卡热量，那么慢跑\(1\)千米大约会消耗（\hspace{1.8em}）千卡热量。妈妈两周共跑了\(100\)千米，大约能消耗（\hspace{1.8em}）千卡热量。\par
\begin{BookAnswer}
先求1千米消耗：\(631\div10=63.1\)千卡。跑100千米消耗\(63.1\times100=6310\)千卡。填63.1、6310。乘法顺序是每千米消耗量乘千米数。
\end{BookAnswer}

% p105_s05
\BookQuestion{一、5}
5. 一个等腰三角形样式的风筝，底角是\(42^\circ\)，它的顶角是（\hspace{1.8em}）°，如果按角分，这是个（\hspace{1.8em}）三角形。\par
\begin{BookAnswer}
等腰三角形两个底角相等，都为42度，顶角为\[\begin{aligned} &180^\circ-42^\circ\times2\\ ={}&180^\circ-84^\circ\\ ={}&96^\circ\end{aligned}\]\(96^\circ>90^\circ\)，所以是钝角三角形；填96、钝角。
\end{BookAnswer}

% p105_s06
\BookQuestion{一、6}
6. 有\(4\)根小棒（如右图），从中挑选\(3\)根，围成一个三角形，这个三角形的周长是（\hspace{1.8em}）厘米。\par
\BookFigure{figures/p105_s06-01.tex}
\begin{BookAnswer}
取三根小棒，要满足“两短边之和大于最长边”。\(2+5=7\)不能围成；\allowbreak{}\(2+5<10\)、\(2+7<10\)也不能围成。只有5、7、10厘米可行，因为\(5+7>10\)。周长\(5+7+10=22\)厘米，填22。
\end{BookAnswer}

% p105_s07
\BookQuestion{一、7}
7. 一个两位小数保留一位小数后约是\(3.8\)，这个小数最大是（\hspace{1.8em}），最小是（\hspace{1.8em}）。\par
\begin{BookAnswer}
保留一位小数后为3.8，原数须满足\(3.75\le x<3.85\)。限定是两位小数，所以最小3.75，最大3.84；3.85已近似为3.9。按空的顺序填3.84、3.75。
\end{BookAnswer}

% p105_s08
\BookQuestion{一、8}
8. 四年级一班有\(42\)名同学，其中会下象棋的有\(19\)人，会下围棋的有\(16\)人，两种棋都不会下的有\(10\)人。两种棋都会下的有（\hspace{1.8em}）人，只会下象棋不会下围棋的有（\hspace{1.8em}）人。\par
\begin{BookAnswer}
至少会一种棋的有\(42-10=32\)人。会象棋19人与会围棋16人相加时，两种都会的人被计了两次，故两种都会：\[\begin{aligned} &19+16-(42-10)\\ ={}&35-32\\ ={}&3\text{(人)}\end{aligned}\]只会象棋：\[\begin{aligned} &19-3\\ ={}&16\text{(人)}\end{aligned}\]依次填3、16。
\end{BookAnswer}

\Needspace{8\baselineskip}

\pdfbookmark[2]{原书第106页}{orig-106}

\subsection*{原书第106页}

% p106_s01
\BookQuestion{二、1}
二、选择（将正确答案的序号填在括号里）。（每题2分，共12分）\par 1. 王大爷要给小菜园围上一圈篱笆，（\hspace{1.8em}）的围法更牢固些。\par
\noindent A\par
\noindent B\par
\noindent C\par
\noindent D\par
\BookFigure{figures/p106_s01-01.tex}
\begin{BookAnswer}
选B。B的篱笆中形成了三角形支撑，利用三角形稳定性，形状不易改变；只由四边形组成的格架容易沿接点变形。
\end{BookAnswer}

% p106_s02
\BookQuestion{二、2}
2. 由我国自主研制的大型多用途民用直升机“吉祥鸟”首飞成功，其最大外吊挂起飞质量为\(13.8\)吨，其中的“8”表示（\hspace{1.8em}）。\par
\noindent A. 8kg\par
\noindent B. 80kg\par
\noindent C. 800kg\par
\noindent D. 8000kg\par
\begin{BookAnswer}
选C。13.8吨中的8在十分位，表示\(0.8\)吨；按进率换算：\[0.8\text{吨}=(0.8\times1000)\text{千克}=800\text{千克}\]所以表示800千克，不是8千克或80千克。
\end{BookAnswer}

% p106_s03
\BookQuestion{二、3}
3. 下面算式中，计算正确的是（\hspace{1.8em}）。\par
\noindent A. \(395-78+22=395-(78+22)\)\par
\noindent B. \(25\times44\times7=25\times4+40\times7\)\par
\noindent C. \(18\times16\div18\times6=(18\times16)\div(18\times6)\)\par
\noindent D. \(99\times48=100\times48-48\)\par
\begin{BookAnswer}
选D。\(99\times48=(100-1)\times48=100\times48-48\)，用的是乘法分配律。\par A把加22误作减22；B没有正确分解44再分配；C同级乘除应从左到右，原左边等于96，而右边把分母改成\(18\times6\)，结果不同。
\end{BookAnswer}

% p106_s04
\BookQuestion{二、4}
4. 把一根细铁丝剪两刀分成三段，围成一个三角形（接头处忽略不计），这两刀可以剪在下图中（\hspace{1.8em}）处。\par
\noindent A. ①②\par
\noindent B. ②③\par
\noindent C. ②④\par
\noindent D. ③④\par
\BookFigure{figures/p106_s04-01.tex}
\begin{BookAnswer}
选C，即②④。把整根铁丝按刻度看成12份，②在第4份末，④在第7份末，两刀后长度为4份、3份、5份，\allowbreak{}\(3+4>5\)，可围成三角形。其他组合的最短两段之和不大于最长段，不能围成。
\end{BookAnswer}

% p106_s05
\BookQuestion{二、5}
5. 下列图示中，不能表示\(0.04\)的是（\hspace{1.8em}）。\par
\noindent A. 计数器\par
\noindent B. 百分格\par
\noindent C. 正方体\par
\noindent D. 数轴\par
\BookFigure{figures/p106_s05-01.tex}
\BookFigure{figures/p106_s05-02.tex}
\BookFigure{figures/p106_s05-03.tex}
\BookFigure{figures/p106_s05-04.tex}
\begin{BookAnswer}
选C。A算盘百分位上4个珠表示4个百分之一；B方格涂4个小格，占\(4\div100=0.04\)；D从0到0.1分成5等份，每份0.02，点在第二格即0.04。C的大正方体按层看，涂色4层占10层中的4层，即\(0.4\)，不是\(0.04\)。
\end{BookAnswer}

% p106_s06
\BookQuestion{二、6}
6. 计算图①和图②中小正方形的总个数时，下面图示的计算过程可以解释（\hspace{1.8em}）。\par
\noindent A. 加法交换律\par
\noindent B. 乘法交换律\par
\noindent C. 乘法结合律\par
\noindent D. 乘法分配律\par
\BookFigure{figures/p106_s06-01.tex}
\begin{BookAnswer}
选D，乘法分配律。图①可按“每行3个、共6行”计数，列\(3\times6\)；图②按“每列3个、共4列”计数，列\(3\times4\)。图①旋转后变成每列3个、共6列，再与图②拼成同高3的长方形；拼图共有\(6+4\)列，按每列3个乘列数，得\(3\times(6+4)\)。\[3\times6+3\times4=3\times(6+4).\]这里两次列式都先写每份的小正方形个数，再写份数；不能把原图①错误描述成“每列3个”。
\end{BookAnswer}

\Needspace{8\baselineskip}

\pdfbookmark[2]{原书第107页}{orig-107}

\subsection*{原书第107页}

% p107_s01
\BookQuestion{三、1}
三、计算。（共30分）\par 1. 直接写得数。（8分）\par \(3.7+4.3=\quad8.36-5.3=\quad0.49\times100=\quad280\div5\div2=\)\par \(9\div100=\quad0.2\times1000=\quad6.4+4.6=\quad35\times25\times4=\)\par
\begin{BookAnswer}
按原题行列，第一行\(8,3.06,49,28\)；第二行\(0.09,200,11,3500\)。\par \(280\div5\div2=280\div(5\times2)=28\)；\allowbreak{}\(35\times25\times4=35\times100=3500\)。\(0.49\times100\)小数点右移两位，\allowbreak{}\(9\div100\)左移两位。
\end{BookAnswer}

% p107_s02
\BookQuestion{三、2（1）}
2. 列竖式计算，带☆的要验算。（10分）\par \(16.8+2.27=\textcolor{answerred}{19.07}\)\par
\begin{BookAnswer}
小数点对齐，16.8补为16.80：\[\begin{array}{r}16.80\\+2.27\\\hline19.07\end{array}\]十分位满十要向个位进1，结果19.07。
\end{BookAnswer}

% p107_s03
\BookQuestion{三、2（2）}
☆\(12-4.5=\textcolor{answerred}{7.5}\)（要验算）\par
\begin{BookAnswer}
12补为12.0，按位退位相减：\[\begin{array}{r}12.0\\-4.5\\\hline7.5\end{array}\]\par 验算：\[\begin{array}{r}7.5\\+4.5\\\hline12.0\end{array}\]
\end{BookAnswer}

% p107_s04
\BookQuestion{三、2（3）}
\(10.02-1.9=\textcolor{answerred}{8.12}\)\par
\begin{BookAnswer}
1.9补为1.90，小数点对齐：\[\begin{array}{r}10.02\\-1.90\\\hline8.12\end{array}\]百分位\(2-0=2\)。十分位不够减，个位又是0，先从十位退1，换成10个一；再从个位退1，换成10个十分之一。\(10-9=1\)，十分位写1；个位剩9，\allowbreak{}\(9-1=8\)。所以得8.12，不能把10.02看成10.20。
\end{BookAnswer}

% p107_s05
\BookQuestion{三、3（1）}
3. 计算下面各题，能简算的要简算。（12分）\par \(180+19\times(106+74)\)\par
\begin{BookAnswer}
先合并括号，再提取180。\[\begin{aligned} &180+19\times(106+74)\\ ={}&180+19\times180\\ ={}&(1+19)\times180\\ ={}&20\times180\\ ={}&3600\end{aligned}\]
\end{BookAnswer}

% p107_s06
\BookQuestion{三、3（2）}
\(40.65-(7.38+10.65)\)\par
\begin{BookAnswer}
先减去容易凑整的10.65。\[\begin{aligned} &40.65-(7.38+10.65)\\ ={}&40.65-10.65-7.38\\ ={}&30-7.38\\ ={}&22.62\end{aligned}\]
\end{BookAnswer}

% p107_s07
\BookQuestion{三、3（3）}
\(25\times201-25\)\par
\begin{BookAnswer}
把减去的25看作一份25。\[\begin{aligned} &25\times201-25\\ ={}&25\times(201-1)\\ ={}&25\times200\\ ={}&5000\end{aligned}\]
\end{BookAnswer}

% p107_s08
\BookQuestion{三、3（4）}
\(840\div[15\times(81-77)]\)\par
\begin{BookAnswer}
先算最内层差，再把15与4凑成整十除数60。\[\begin{aligned}&840\div[15\times(81-77)]\\={}&840\div[15\times4]\\={}&840\div60\\={}&14\end{aligned}\]
\end{BookAnswer}

% p107_s09
\BookQuestion{四、（1）}
四、操作。（共6分）\par 下图是一块平行四边形露营地，奇思在这块露营地的小路上跑步，淘气在草坪上跑步。\par （1）淘气跑的路线与小路平行，请你画一条直线表示淘气跑步的路线。\par
\BookFigure{figures/p107_s09-01.tex}
\begin{BookAnswer}
过淘气所在的黑点作一条与小路边界平行的直线。应比较整条路线的方向，不能只凭目测画成水平线；左图红线已通过淘气所在点，并与斜放的小路平行。
\BookFigure{figures/p107_s09-02.tex}
\end{BookAnswer}

% p107_s10
\BookQuestion{四、（2）}
（2）笑笑想走到奇思所在的这条小路上，请你画出她去这条小路的最短路线。\par
\BookFigure{figures/p107_s10-01.tex}
\begin{BookAnswer}
从笑笑所在的右下角点，向奇思所在小路的最近边界作垂线段；垂足必须落在小路边界上，直角标记的两边分别沿小路边界和所作路线。点到直线的垂线段最短，所以该路线是到小路的最短路线。
\BookFigure{figures/p107_s10-02.tex}
\end{BookAnswer}

\Needspace{8\baselineskip}

\pdfbookmark[2]{原书第108页}{orig-108}

\subsection*{原书第108页}

% p108_s01
\BookQuestion{五、1}
五、解决问题。（共29分）\par 1. 妈妈在微信群里发了\(50\)元的红包，爸爸抢到\(15.65\)元，小芳抢到\(19.7\)元。小芳比爸爸多抢多少元？妈妈发的红包还剩下多少元？（5分）\par
\begin{BookAnswer}
先求“多抢多少元”，再求“红包剩多少元”。\par \textbf{多抢：直接求两人金额的差}\par 小芳19.7元比爸爸15.65元多多少，用大数减小数。\[\begin{aligned} &19.7-15.65\\ ={}&19.70-15.65\\ ={}&4.05\text{(元)}\end{aligned}\]答：小芳比爸爸多抢4.05元。\par \textbf{剩余方法一：依次扣除}\par 从妈妈发的50元依次扣掉爸爸与小芳抢到的金额。\[\begin{aligned} &50-15.65-19.7\\ ={}&34.35-19.7\\ ={}&14.65\text{(元)}\end{aligned}\]答：红包还剩14.65元。\par \textbf{剩余方法二：先合计两人抢到的钱}\par 先加两人的金额，再从50元中一次扣除。\[\begin{aligned} &50-(15.65+19.7)\\ ={}&50-35.35\\ ={}&14.65\text{(元)}\end{aligned}\]答：红包还剩14.65元。\par 按题面19.7与15.65核算，书后18.65元与原数据不符，不采用。\par 用剩余的钱加上两人已抢到的钱检验：\(14.65+15.65+19.70=30.30+19.70=50.00\)元，正好还原妈妈发出的50元。
\end{BookAnswer}

% p108_s02
\BookQuestion{五、2}
2. 小明在图书室借阅了《科学探索之谜》这本书，如果每天看\(40\)页，\allowbreak{}\(11\)天才能看完。学校要求按时归还，需要提前\(3\)天看完，他平均每天要看多少页？（6分）\par
\begin{BookAnswer}
\textbf{方法一：总页数除以实际天数，简便计算}每天40页乘11天为总页数，括成整体；提前3天，实际8天。利用乘积除以一个数的性质，40可被8整除，先用40除以8再乘11：\[\begin{aligned}&(40\times11)\div(11-3)\\={}&(40\times11)\div8\\={}&(40\div8)\times11\\={}&5\times11\\={}&55\text{(页)}\end{aligned}\]答：平均每天要看55页。\par\textbf{方法二：把少用3天的任务平均补给8天}\[\begin{aligned}&40+(40\times3)\div(11-3)\\={}&40+120\div8\\={}&40+15\\={}&55\text{(页)}\end{aligned}\]答：平均每天要看55页。
\end{BookAnswer}

% p108_s03
\BookQuestion{五、3}
3. 学校为\(45\)名学生订购合唱服，上衣每件\(62\)元，裤子每条\(38\)元，买这批合唱服一共需要多少元？（6分）\par
\begin{BookAnswer}
每名学生各买一件上衣和一条裤子，可先把每套单价62与38凑100，再乘45套。分别列上衣、裤子总价时，用乘法分配律合并共同的45；人数不能因此变成90。\par \[\begin{aligned}&62\times45+38\times45\\={}&(62+38)\times45\\={}&100\times45\\={}&4500\text{(元)}\end{aligned}\]答：买这批合唱服一共需要4500元。
\end{BookAnswer}

% p108_s04
\BookQuestion{五、4}
4. 修路队要抢修一段公路。前\(4\)天共修路\(2500\)米，后\(3\)天平均每天修路\(800\)米，这支修路队这一周平均每天修路多少米？（6分）\par
\begin{BookAnswer}
\textbf{方法一：总路程除以总天数}\par 后3天平均每天800米，先求后3天总量，再加前4天的2500米。\[\begin{aligned} &(2500+800\times3)\div(4+3)\\ ={}&(2500+2400)\div7\\ ={}&4900\div7\\ ={}&700\text{(米)}\end{aligned}\]答：这一周平均每天修路700米。\par \textbf{方法二：以每天800米为基准}\par 前4天若每天修800米，应修3200米，比实际多700米；把这700米平均扣给7天。\[\begin{aligned} &800-(800\times4-2500)\div7\\ ={}&800-(3200-2500)\div7\\ ={}&800-700\div7\\ ={}&800-100\\ ={}&700\text{(米)}\end{aligned}\]答：这一周平均每天修路700米。\par
\end{BookAnswer}

% p108_s05
\BookQuestion{五、5（统计表与填空）}
5. 下面是某小学教师的年龄记录。（6分）\par \(37\quad28\quad26\quad23\quad32\quad39\quad28\quad34\quad45\quad41\quad26\quad46\quad37\quad25\)\par \(57\quad27\quad55\quad33\quad34\quad27\quad52\quad42\quad24\quad50\quad48\quad34\quad40\quad51\)\par 根据以上数据填写下面的统计表。\par 教师年龄在（\hspace{1.8em}）范围内的人数最多，（\hspace{1.8em}）范围内的人数最少。\par
\par\smallskip\Needspace{4\baselineskip}\begingroup\small\setlength{\tabcolsep}{3pt}
\noindent\begin{tabularx}{\linewidth}{|*{5}{>{\centering\arraybackslash}X|}}\hline
年龄 & 20～29岁 & 30～39岁 & 40～49岁 & 50～59岁\\\hline
人数 &  &  &  & \\\hline
\end{tabularx}\par
\endgroup\smallskip
\begin{BookAnswer}
逐人按年龄段登记，不能把端点跨组：\par \[\begin{array}{c|rrrr}\text{年龄/岁}&20\sim29&30\sim39&40\sim49&50\sim59\\\hline\text{人数}&9&8&6&5\end{array}\]\par 总人数核对\(9+8+6+5=28\)。20至29岁9人最多，50至59岁5人最少，依次填“20～29岁”“50～59岁”。\par 可核对分组：20～29岁有28、26、23、28、26、25、27、27、24；30～39岁有37、32、39、34、37、33、34、34；40～49岁有45、41、46、42、48、40；50～59岁有57、55、52、50、51。
\end{BookAnswer}
